\documentclass[12pt,a4paper]{ctexart}
\usepackage{amsmath,amssymb,bm}
\usepackage{geometry}
\usepackage{enumitem}
\usepackage{booktabs}
\usepackage{hyperref}

\geometry{left=2.2cm,right=2.2cm,top=2.2cm,bottom=2.2cm}
\setlist[enumerate]{itemsep=0.7em,topsep=0.5em}

\newcommand{\Rey}{\mathrm{Re}}
\newcommand{\dd}{\mathrm d}

\title{第八章 粘性流动 作业-1：完整答案}
\author{}
\date{}

\begin{document}
\maketitle

\section*{1. 判断层流或湍流}

圆管流动雷诺数为
\[
\Rey=\frac{Vd}{\nu}.
\]

\paragraph{(1)}
\[
\Rey=\frac{1\times0.2}{1\times10^{-6}}
=2.0\times10^5.
\]
远大于临界值，故为湍流：
\[
\boxed{\text{湍流}}.
\]

\paragraph{(2)}
\[
\Rey=\frac{0.2\times0.15}{2\times10^{-5}}
=1.5\times10^3.
\]
小于 \(2000\)，故为层流：
\[
\boxed{\text{层流}}.
\]

\section*{3. 书 8.4}

已知
\[
d=75\mathrm{mm}=0.075\mathrm m,\quad
Q=7\mathrm{L/s}=0.007\mathrm{m^3/s},\quad
\rho=800\mathrm{kg/m^3},\quad
\tau_w=48\mathrm{N/m^2}.
\]

平均速度
\[
\bar V=\frac{Q}{A}
=\frac{Q}{\pi d^2/4}
=\frac{0.007}{\pi(0.075)^2/4}
\approx1.584\mathrm{m/s}.
\]

圆管层流中
\[
\tau_w=\mu\left|\frac{\dd u}{\dd r}\right|_{r=R}
=\frac{8\mu \bar V}{d}.
\]

所以
\[
\mu=\frac{\tau_w d}{8\bar V}
=\frac{48\times0.075}{8\times1.584}
\approx0.284\mathrm{kg/(m\cdot s)}.
\]

验算雷诺数：
\[
\Rey=\frac{\rho \bar V d}{\mu}
=\frac{800\times1.584\times0.075}{0.284}
\approx335<2000,
\]
层流假设成立。

\[
\boxed{\mu\approx0.28\mathrm{kg/(m\cdot s)}}.
\]

\section*{4. 书 8.5}

钢球半径
\[
a=\frac{1.5}{2}\mathrm{mm}=0.75\times10^{-3}\mathrm m.
\]

钢球等速下降速度
\[
U=\frac{0.500}{56}
\approx8.93\times10^{-3}\mathrm{m/s}.
\]

油的密度
\[
\rho=0.95\times1000=950\mathrm{kg/m^3}.
\]

小球体积
\[
V_s=\frac43\pi a^3.
\]

小雷诺数下，阻力服从斯托克斯公式：
\[
D=6\pi\mu aU.
\]

等速下落时合力为零：
\[
mg-\rho gV_s=6\pi\mu aU.
\]

因此
\[
\mu=\frac{mg-\rho gV_s}{6\pi aU}.
\]

代入
\[
m=13.7\times10^{-6}\mathrm{kg},\quad
a=0.75\times10^{-3}\mathrm m,\quad
U=8.93\times10^{-3}\mathrm{m/s},
\]
得
\[
\mu\approx0.934\mathrm{kg/(m\cdot s)}.
\]

验算：
\[
\Rey=\frac{\rho Ud}{\mu}
=\frac{950\times8.93\times10^{-3}\times1.5\times10^{-3}}{0.934}
\approx0.014<1.
\]
故小雷诺数假设成立。

\[
\boxed{\mu\approx0.934\mathrm{kg/(m\cdot s)}}.
\]

\section*{5. 书 8.6}

取 \(y=0\) 为下板，\(y=h_1\) 为两流体交界面，\(y=h_1+h_2\) 为上板。
下层为流体 1，上层为流体 2。

由于全场压强为常数，且不计质量力，动量方程化为
\[
\mu_i\frac{\dd^2u_i}{\dd y^2}=0.
\]

因此两层速度均为线性分布：
\[
u_1=A_1y+B_1,\qquad 0\le y\le h_1,
\]
\[
u_2=A_2y+B_2,\qquad h_1\le y\le h_1+h_2.
\]

边界条件为
\[
u_1(0)=0,
\]
\[
u_2(h_1+h_2)=U,
\]
\[
u_1(h_1)=u_2(h_1),
\]
\[
\mu_1\frac{\dd u_1}{\dd y}
=
\mu_2\frac{\dd u_2}{\dd y}.
\]

设两层剪应力相同，为
\[
\tau=\mu_1\frac{\dd u_1}{\dd y}
=\mu_2\frac{\dd u_2}{\dd y}.
\]

总速度差为
\[
U=\frac{\tau h_1}{\mu_1}+\frac{\tau h_2}{\mu_2}.
\]

所以
\[
\tau=
\frac{U}{h_1/\mu_1+h_2/\mu_2}
=\frac{U\mu_1\mu_2}{\mu_2h_1+\mu_1h_2}.
\]

下层速度：
\[
u_1(y)=\frac{\tau}{\mu_1}y
=
\frac{U\mu_2 y}{\mu_2h_1+\mu_1h_2},
\qquad 0\le y\le h_1.
\]

上层速度：
\[
u_2(y)=u_1(h_1)+\frac{\tau}{\mu_2}(y-h_1),
\]
即
\[
u_2(y)
=
\frac{U\left[\mu_2h_1+\mu_1(y-h_1)\right]}
{\mu_2h_1+\mu_1h_2},
\qquad h_1\le y\le h_1+h_2.
\]

剪应力分布为常数：
\[
\boxed{
\tau_{xy}=
\frac{U\mu_1\mu_2}{\mu_2h_1+\mu_1h_2}
}.
\]

答案：
\[
\boxed{
u_1(y)=
\frac{U\mu_2 y}{\mu_2h_1+\mu_1h_2},
\quad 0\le y\le h_1
}
\]
\[
\boxed{
u_2(y)=
\frac{U\left[\mu_2h_1+\mu_1(y-h_1)\right]}
{\mu_2h_1+\mu_1h_2},
\quad h_1\le y\le h_1+h_2
}
\]
\[
\boxed{
\tau_{xy}=
\frac{U\mu_1\mu_2}{\mu_2h_1+\mu_1h_2}
}.
\]

\section*{6. 书 8.10}

小球半径为 \(a\)，密度为 \(\rho_s\)，流体密度为 \(\rho\)，流体粘度为 \(\mu\)。

小球重力：
\[
G=\frac43\pi a^3\rho_s g.
\]

浮力：
\[
B=\frac43\pi a^3\rho g.
\]

小雷诺数下，斯托克斯阻力为
\[
D=6\pi\mu aU.
\]

小球达到最终速度 \(U_t\) 时，受力平衡：
\[
G-B-D=0.
\]

即
\[
\frac43\pi a^3(\rho_s-\rho)g
=6\pi\mu aU_t.
\]

解得
\[
U_t=
\frac{\frac43\pi a^3(\rho_s-\rho)g}{6\pi\mu a}
=
\frac{2a^2(\rho_s-\rho)g}{9\mu}.
\]

故
\[
\boxed{
U_t=\frac{2a^2(\rho_s-\rho)g}{9\mu}
}.
\]

\section*{7. 管内层流：不同方向的压降和雷诺数}

已知
\[
\rho=1.24\mathrm{kg/m^3},\quad
\mu=0.001\mathrm{kg/(m\cdot s)},\quad
D=0.08\mathrm m,\quad
\tau_w=0.72\mathrm{N/m^2}.
\]

圆管层流中
\[
\tau_w=\frac{8\mu \bar V}{D}.
\]

所以
\[
\bar V=\frac{\tau_wD}{8\mu}
=\frac{0.72\times0.08}{8\times0.001}
=7.2\mathrm{m/s}.
\]

雷诺数
\[
\Rey=\frac{\rho \bar VD}{\mu}
=\frac{1.24\times7.2\times0.08}{0.001}
\approx714.
\]

故三种情况均为层流。

管流中有
\[
\tau_w=\frac{D}{4}\left(-\frac{\dd p}{\dd x}+\rho g_x\right),
\]
所以
\[
\frac{\dd p}{\dd x}
=\rho g_x-\frac{4\tau_w}{D}.
\]

其中
\[
\frac{4\tau_w}{D}
=\frac{4\times0.72}{0.08}
=36\mathrm{Pa/m}.
\]

\paragraph{(1) 水平管}
\[
g_x=0.
\]
\[
\frac{\dd p}{\dd x}=-36\mathrm{Pa/m}.
\]
即沿程压降
\[
-\frac{\dd p}{\dd x}=36\mathrm{Pa/m}.
\]

\paragraph{(2) 向上垂直流动}
取 \(x\) 沿流动向上，则
\[
g_x=-g.
\]
\[
\frac{\dd p}{\dd x}
=-\rho g-36
=-1.24\times9.81-36
\approx-48.2\mathrm{Pa/m}.
\]
即沿程压降
\[
-\frac{\dd p}{\dd x}\approx48.2\mathrm{Pa/m}.
\]

\paragraph{(3) 向下垂直流动}
取 \(x\) 沿流动向下，则
\[
g_x=+g.
\]
\[
\frac{\dd p}{\dd x}
=\rho g-36
=1.24\times9.81-36
\approx-23.8\mathrm{Pa/m}.
\]
即沿程压降
\[
-\frac{\dd p}{\dd x}\approx23.8\mathrm{Pa/m}.
\]

答案汇总：
\[
\boxed{
\Rey\approx714
}
\]
\[
\boxed{
\begin{array}{c|c|c}
\text{情况} & dp/dx\ (\mathrm{Pa/m}) & -dp/dx\ (\mathrm{Pa/m})\\
\hline
\text{水平} & -36.0 & 36.0\\
\text{向上} & -48.2 & 48.2\\
\text{向下} & -23.8 & 23.8
\end{array}
}
\]

\section*{8. 倾斜平板间流动}

取 \(x\) 沿斜面向下，\(y\) 从下板指向上板。
下板为
\[
y=0,\qquad u(0)=0.
\]
上板为
\[
y=H,\qquad u(H)=-U_0,
\]
因为上板向上拉动，与 \(x\) 正方向相反。

沿流动方向无压力差，动量方程为
\[
\mu\frac{\dd^2u}{\dd y^2}+\rho g\sin\theta=0.
\]

令
\[
A=\frac{\rho g\sin\theta}{\mu}.
\]
则
\[
\frac{\dd^2u}{\dd y^2}=-A.
\]

积分得
\[
u=-\frac A2y^2+C_1y+C_2.
\]

由 \(u(0)=0\)，得 \(C_2=0\)。

又已知
\[
u\left(\frac H2\right)=0.
\]

所以
\[
0=-\frac A2\left(\frac H2\right)^2+C_1\frac H2,
\]
得
\[
C_1=\frac{AH}{4}.
\]

因此
\[
u=-\frac A2y^2+\frac{AH}{4}y.
\]

代入上壁面条件
\[
u(H)=-U_0,
\]
得
\[
-U_0=-\frac A2H^2+\frac A4H^2=-\frac{AH^2}{4}.
\]

故
\[
U_0=\frac{AH^2}{4}
=\frac{\rho g\sin\theta\,H^2}{4\mu}.
\]

代入
\[
\rho=1.25,\quad
H=0.03\mathrm m,\quad
\theta=15^\circ,\quad
\mu=0.0001,
\]
得
\[
U_0=
\frac{1.25\times9.81\times\sin15^\circ\times(0.03)^2}{4\times0.0001}
\approx7.14\mathrm{m/s}.
\]

剪应力
\[
\tau=\mu\frac{\dd u}{\dd y}
=\mu\left(-Ay+\frac{AH}{4}\right)
=\rho g\sin\theta\left(\frac H4-y\right).
\]

下壁面：
\[
\tau_0=\rho g\sin\theta\frac H4.
\]

上壁面：
\[
\tau_H=\rho g\sin\theta\left(\frac H4-H\right)
=-\frac{3\rho gH\sin\theta}{4}.
\]

故切应力大小为
\[
|\tau_0|
=\frac{\rho gH\sin\theta}{4}
\approx2.38\times10^{-2}\mathrm{Pa},
\]
\[
|\tau_H|
=\frac{3\rho gH\sin\theta}{4}
\approx7.14\times10^{-2}\mathrm{Pa}.
\]

答案：
\[
\boxed{U_0\approx7.14\mathrm{m/s}}
\]
\[
\boxed{|\tau_{\text{下壁}}|\approx2.38\times10^{-2}\mathrm{Pa}}
\]
\[
\boxed{|\tau_{\text{上壁}}|\approx7.14\times10^{-2}\mathrm{Pa}}
\]

\section*{9. 证明静止物面上 \(\bm\Omega\cdot\bm n=0\)}

在静止固体壁面上，粘性流体满足无滑移条件。
取壁面局部坐标，使 \(x,y\) 为壁面切向坐标，\(z\) 为法向坐标。
则壁面为 \(z=0\)，法向量为
\[
\bm n=\bm k.
\]

无滑移条件给出
\[
u(x,y,0)=0,\qquad v(x,y,0)=0,\qquad w(x,y,0)=0.
\]

涡量的法向分量为
\[
\bm\Omega\cdot\bm n
=\Omega_z
=\frac{\partial v}{\partial x}
-\frac{\partial u}{\partial y}.
\]

因为在壁面上 \(u\) 和 \(v\) 沿整个壁面均为零，
所以它们沿壁面切向的导数也为零：
\[
\left.\frac{\partial v}{\partial x}\right|_{z=0}=0,
\qquad
\left.\frac{\partial u}{\partial y}\right|_{z=0}=0.
\]

因此
\[
\left.\Omega_z\right|_{z=0}=0.
\]

即
\[
\boxed{
\bm\Omega\cdot\bm n=0.
}
\]

物理解释：静止壁面上任意贴壁小闭合曲线的速度环量为零，因此穿过该小面积的涡通量也为零。

\section*{10. 证明 \(\nabla^2\psi=-\Omega\)}

二维不可压流动中，引入流函数 \(\psi\)，取约定
\[
u=\frac{\partial\psi}{\partial y},
\qquad
v=-\frac{\partial\psi}{\partial x}.
\]

涡量只有 \(z\) 方向分量：
\[
\bm\Omega=\Omega(x,y)\bm k.
\]

其中
\[
\Omega=\frac{\partial v}{\partial x}
-\frac{\partial u}{\partial y}.
\]

代入流函数表达式：
\[
\frac{\partial v}{\partial x}
=\frac{\partial}{\partial x}
\left(-\frac{\partial\psi}{\partial x}\right)
=-\frac{\partial^2\psi}{\partial x^2},
\]
\[
\frac{\partial u}{\partial y}
=\frac{\partial}{\partial y}
\left(\frac{\partial\psi}{\partial y}\right)
=\frac{\partial^2\psi}{\partial y^2}.
\]

所以
\[
\Omega
=-\frac{\partial^2\psi}{\partial x^2}
-\frac{\partial^2\psi}{\partial y^2}
=-\nabla^2\psi.
\]

因此
\[
\boxed{
\nabla^2\psi=-\Omega.
}
\]

\end{document}
